\begin{figure}[t]\centering
% ============================ TUNING KNOBS ============================
\def\rowp{13pt}      % baseline-to-baseline pitch of the program rows
\def\cop{4}          % column (in characters) where the opcode starts
\def\cdst{12.5}      % column of the destination operand
\def\carr{17.4}      % column of the <- arrow
\def\csrc{20}        % column of the first source operand
\def\crail{31.6}     % column of the long rf rail (right of every row)
\def\cflr{34.4}      % column of the r flag
\def\cflc{36.4}      % column of the c flag
\def\cbr{-0.9}       % column of the "added" bracket
\def\domeA{9pt}      % control-point height of the ddi arcs in I1 and I3 (room above: short branch rows)
\def\domeB{4.4pt}    % control-point height of the ddi arc in I4 (stays under the incoming rf)
\def\gdrop{1.2pt}    % how far an rf line drops below its source operand before turning
\def\keygap{2pt}     % gap between the key line and the first row
% ======================================================================


% width of one typewriter character at the body size; every column is a multiple of it
\setbox0=\hbox{\footnotesize\ttfamily 0}\edef\cw{\the\wd0}
\begin{tikzpicture}[
  >=latex,
  every node/.style={font=\footnotesize, inner sep=0pt,
    execute at begin node={\hyphenpenalty=10000\exhyphenpenalty=10000}},
  tok/.style={font=\footnotesize\ttfamily, anchor=base west},
  mem/.style={tok, draw, inner sep=0.9pt},
  lab/.style={font=\scriptsize\itshape},
  flag/.style={anchor=base},
]
\newcommand{\tk}[4]{\node[tok] (#1) at ({(#2)*\cw},{-(#3)*\rowp}) {#4};}

% ------------------------------ program ------------------------------
% I0 (added)
\tk{l0}{0}{0}{I0:}  \tk{o0}{\cop}{0}{branch\_n}
% I1 (added)
\tk{l1}{0}{1}{I1:}  \tk{o1}{\cop}{1}{load}  \tk{d1}{\cdst}{1}{reg0}
\node[tok] at ({\carr*\cw},{-1*\rowp}) {$\leftarrow$};
\node[mem] (m1) at ({\csrc*\cw},{-1*\rowp}) {mem 0};
% I2
\tk{l2}{0}{2}{I2:}  \tk{o2}{\cop}{2}{branch\_n}
% I3
\tk{l3}{0}{3}{I3:}  \tk{o3}{\cop}{3}{load}  \tk{d3}{\cdst}{3}{reg1}
\node[tok] at ({\carr*\cw},{-3*\rowp}) {$\leftarrow$};
\node[mem] (m3) at ({\csrc*\cw},{-3*\rowp}) {mem 1};
% I4
\tk{l4}{0}{4}{I4:}  \tk{o4}{\cop}{4}{other\_n}  \tk{d4}{\cdst}{4}{reg2}
\node[tok] at ({\carr*\cw},{-4*\rowp}) {$\leftarrow$};
\tk{s4a}{\csrc}{4}{reg1}  \tk{s4c}{\csrc+4}{4}{,}  \tk{s4b}{\csrc+6}{4}{reg0}
% I5
\tk{l5}{0}{5}{I5:}  \tk{o5}{\cop}{5}{load}  \tk{d5}{\cdst}{5}{?}
\node[tok] at ({\carr*\cw},{-5*\rowp}) {$\leftarrow$};
\tk{s5o}{\csrc}{5}{[}  \tk{s5}{\csrc+1}{5}{reg2}  \tk{s5c}{\csrc+5}{5}{]}

% ------------------------- resolved / committed -------------------------
\foreach \row/\r/\c in {0/r/c, 1/r/c, 2//, 3/r/, 4/r/, 5/r/} {
  \node[flag] at ({\cflr*\cw+0.5*\cw},{-\row*\rowp}) {\r};
  \node[flag] at ({\cflc*\cw+0.5*\cw},{-\row*\rowp}) {\c};
}
\node[lab, anchor=base east] at ({\cflc*\cw+1*\cw},{\rowp-\keygap}) {r = resolved,\ c = committed};

% ------------------------------- edges -------------------------------
% Convention: an arc above a row is a ddi edge (inside one instruction); a routed line
% between rows is an rf edge (output operand of one instruction -> input operand of a later one).
% ddi inside I1 and I3 (inmem -> outreg)
\draw[->] (m1.north) .. controls +(0,\domeA) and +(0,\domeA) .. (d1.north)
  node[lab, pos=0.5, above=0.6pt] {ddi};
\draw[->] (m3.north) .. controls +(0,\domeA) and +(0,\domeA) .. (d3.north)
  node[lab, pos=0.5, above=0.6pt] {ddi};
% ddi inside I4 (reg1 -> reg2): a low arc that stays under the rf landing on reg1
\draw[->] ($(s4a.north west)+(0.6*\cw,0)$) .. controls +(0,\domeB) and +(0,\domeB) .. ($(d4.north east)+(-0.6*\cw,0)$);
% rf I3.reg1 -> I4.reg1
\draw[->] (d3.south) -- ++(0,-\gdrop) -| ($(s4a.north)+(0.7*\cw,0)$);
% rf I4.reg2 -> I5.reg2
\draw[->] (d4.south) -- ++(0,-\gdrop) -| (s5.north);
% rf I1.reg0 -> I4.reg0: under I1, then down the rail right of every row
\coordinate (rail) at ({\crail*\cw},0);
\draw[->] (d1.south) -- ++(0,-\gdrop) coordinate (g1) -- (g1 -| rail) -- ($(rail |- s4b.center)$) -- (s4b.east);
\node[lab, anchor=west] at ($(rail |- l3.base)+(1.2pt,0)$) {rf};

% ------------------------- added-instruction bracket -------------------------
\coordinate (bt) at ({\cbr*\cw},{0.72*\rowp});
\coordinate (bb) at ({\cbr*\cw},{-1*\rowp-0.3*\rowp});
\draw ($(bt)+(0.45*\cw,0)$) -- (bt) -- (bb) -- ($(bb)+(0.45*\cw,0)$);
\node[lab, anchor=east, align=right] at ($(bt)!0.5!(bb)+(-1.5pt,0)$) {added by\\interleaving};

\end{tikzpicture}
\end{figure}
